%
% bootstrapping-pharo.tex
%
% Bootstrapping a Smalltalk Messaging SDK from an External LLM Agent
%
% Compile: pdflatex bootstrapping-pharo.tex
%

\documentclass[a4paper,10pt,fleqn]{article}
\usepackage[margin=1in]{geometry}
\usepackage[colorlinks=true,linkcolor=blue,citecolor=blue,urlcolor=blue]{hyperref}
\usepackage{booktabs}
\usepackage{array}
\usepackage{tabularx}
\usepackage{longtable}
\usepackage{enumitem}
\usepackage{listings}
\usepackage{xcolor}

\newcolumntype{L}[1]{>{\raggedright\arraybackslash}p{#1}}
\newcolumntype{Y}{>{\raggedright\arraybackslash}X}

\lstset{
  basicstyle=\ttfamily\small,
  keywordstyle=\color{blue!80!black},
  commentstyle=\color{green!50!black},
  stringstyle=\color{red!60!black},
  breaklines=true,
  frame=single,
  framerule=0.3pt,
  rulecolor=\color{gray!50},
  backgroundcolor=\color{gray!5},
  columns=fullflexible,
}

\newcommand{\cls}[1]{\texttt{#1}}
\newcommand{\sel}[1]{\texttt{\##1}}
\newcommand{\pkg}[1]{\texttt{#1}}
\newcommand{\phase}[1]{\textbf{Phase~#1}}

\begin{document}

\title{Bootstrapping a Smalltalk Messaging SDK\\from an External LLM Agent}
\author{Punt Labs}
\date{April 2026}
\hypersetup{
  pdftitle={Bootstrapping a Smalltalk Messaging SDK from an External LLM Agent},
  pdfauthor={Punt Labs},
  pdfsubject={Technical Report},
  pdfcreator={pdflatex}
}
\maketitle

\tableofcontents
\newpage

% ===================================================================
\section{Introduction}
% ===================================================================

This document describes the bootstrapping process by which an LLM agent
running outside a Smalltalk image (Claude Code, a TypeScript-based CLI
tool) progressively builds a Claude Messaging SDK \emph{inside} a Pharo
Smalltalk image.  The bootstrap produced a fully-designed remote driver
(\href{https://github.com/punt-labs/postern}{Postern}) and the
Messaging SDK that ships from this repository.  Agent self-improvement
--- a second LLM agent running inside the image with its own tools ---
continues in the separate \texttt{claude-agent-sdk-pharo} repository.

The process has four phases, each transferring more capability into the
image until the external agent is no longer needed for Smalltalk
development.  The key technical innovation is the HTTP eval server: a
2-line Zinc endpoint that bridges the gap between the filesystem-oriented
external agent and the image-oriented Smalltalk runtime.

\subsection{Context}

Two projects form the planned stack:

\begin{center}
\begin{tabularx}{\linewidth}{@{}llY@{}}
\toprule
\textbf{Project} & \textbf{Status} & \textbf{Role} \\
\midrule
\texttt{anthropic-sdk-pharo} & Shipping (this repository) &
  Claude Messaging SDK for Pharo: API client, message and content-block
  types, streaming, errors, server tools \\
\texttt{claude-agent-sdk-pharo} & Planned &
  Agent runtime layer: system prompt, tool dispatch, conversation
  management, image introspection.  Depends on the Messaging SDK \\
\bottomrule
\end{tabularx}
\end{center}

\subsection{The Problem}

An LLM agent that operates on files (Claude Code) cannot directly
modify a running Smalltalk image.  Smalltalk has no source files in the
traditional sense --- the image \emph{is} the program.  Classes are
created by sending messages to superclasses, methods are compiled by
sending \sel{compile:} to classes, and objects exist in a live heap.

Claude Code's tools (\texttt{Read}, \texttt{Write}, \texttt{Edit},
\texttt{Bash}) operate on the filesystem.  They can write a \texttt{.st}
file, but that file is inert until something loads it into the image.
The question is: how does an external agent bootstrap an internal agent
when the two operate on fundamentally different substrates?

% ===================================================================
\section{The Bootstrap Mechanism}
% ===================================================================

\subsection{Mechanism A: HTTP Eval Server}

The bridge is a Zinc HTTP server running inside the Pharo image that
accepts Smalltalk source code as POST body and evaluates it:

\begin{lstlisting}[language={},title={The minimal bridge (2 lines)}]
ZnServer startDefaultOn: 8422.
ZnServer default delegate: ZnReadEvalPrintDelegate new.
\end{lstlisting}

The external agent interacts with the image via:

\begin{lstlisting}[language={},title={External agent sends Smalltalk}]
curl -s -X POST http://localhost:8422/repl \
  -H "Content-Type: text/plain" \
  -d "3 + 4"
# => 7
\end{lstlisting}

This is not a remote procedure call.  It is arbitrary code execution
in the image's process space.  The external agent can:

\begin{itemize}
  \item Define classes (\texttt{Object subclass: \#Foo ...})
  \item Compile methods (\texttt{Foo compile: 'bar \^{} 42'})
  \item Create objects (\texttt{Foo new})
  \item Query the system (\texttt{Smalltalk globals keys})
  \item Open Morphic windows (\texttt{'hello' asStringMorph openInWorld})
  \item Inspect results (\texttt{Transcript contents})
\end{itemize}

\textbf{Key property:} the eval server runs in a Zinc worker thread.
Evaluated code executes in that thread's context, not on the Morphic UI
thread.  This means blocking network calls (like API requests) do not
freeze the UI.  However, Morphic modifications from evaluated code
require \texttt{WorldState defer:} for thread safety.

This 2-line server is the seed from which
\href{https://github.com/punt-labs/postern}{Postern} evolved into a
fully-designed product.

\subsection{Mechanism B: File-In}

For batch operations (defining multiple classes with correct load
order), the external agent writes Pharo chunk-format \texttt{.st} files
and triggers file-in via the eval server:

\begin{lstlisting}[language={},title={File-in via eval server}]
curl -s -X POST http://localhost:8422/repl \
  -H "Content-Type: text/plain" \
  -d "CodeImporter evaluateFileNamed: '/path/to/MyClass.st'"
\end{lstlisting}

Or headless via the CLI for initial setup:

\begin{lstlisting}[language={},title={Headless file-in}]
pharo --headless Pharo.image eval --save \
  "CodeImporter evaluateFileNamed: 'src/MyClass.st'"
\end{lstlisting}

The Makefile automates this with dependency-ordered loading to avoid
undeclared-variable warnings.

\subsection{Mechanism C: Iceberg (Git)}

Pharo~12 includes Iceberg, a native git client built on libgit2.  Once
the repository uses Tonel source format (one \texttt{.class.st} file per
class, one directory per package), Iceberg provides a third code path:

\begin{itemize}
  \item \textbf{Load from repository.}  Iceberg reads Tonel files from
    \texttt{src/} and loads classes into the image.
  \item \textbf{Commit from image.}  Changes in the System Browser are
    written to Tonel files and committed via libgit2.
  \item \textbf{Push and fetch.}  SSH (ed25519), remote tracking, and
    branch management are built in.
\end{itemize}

Supported operations:

\begin{center}
\begin{tabular}{@{}ll@{}}
\toprule
\textbf{Operation} & \textbf{Supported} \\
\midrule
Open local repo       & Yes \\
Read commit log       & Yes \\
Create/checkout branch & Yes \\
Stage and commit      & Yes \\
Push / fetch          & Yes \\
Merge                 & No (use CLI) \\
Tags                  & Yes \\
\bottomrule
\end{tabular}
\end{center}

\textbf{When to use each mechanism:}

\begin{center}
\begin{tabular}{@{}lll@{}}
\toprule
\textbf{Mechanism} & \textbf{Best for} & \textbf{Limitation} \\
\midrule
A (HTTP eval) & Interactive dev, testing & No version control \\
B (File-in)   & Batch load, rebuild     & No incremental updates \\
C (Iceberg)   & Production workflow     & Requires Tonel format \\
\bottomrule
\end{tabular}
\end{center}

\subsection{Combined Bootstrap Flow}

\begin{enumerate}
  \item \texttt{make setup} downloads a fresh Pharo VM and image.
  \item \texttt{make setup} files in all \texttt{src/*.st} in
    dependency order (headless, saved to image).
  \item \texttt{make start} launches the image with GUI and starts the
    eval server.
  \item The external agent (Claude Code) sends Smalltalk via HTTP POST
    to develop, test, and iterate.
  \item \texttt{make rebuild} proves all code is self-contained by
    rebuilding from a fresh image.
\end{enumerate}

% ===================================================================
\section{The Four Phases}
% ===================================================================

\subsection{Phase 1: External Agent Builds the SDK}

\begin{center}
\begin{tabularx}{\linewidth}{@{}lY@{}}
\toprule
\textbf{Builder} & Claude Code (TypeScript CLI) \\
\textbf{Target} & \texttt{anthropic-sdk-pharo} \\
\textbf{Substrate} & Filesystem (\texttt{.st} files) + eval server \\
\textbf{Tests} & \texttt{curl} against eval server, manual verification \\
\bottomrule
\end{tabularx}
\end{center}

Claude Code writes Smalltalk source files (\texttt{ClaudeClient.st},
\texttt{ClaudeMessage.st}, etc.), files them into the running image via
the eval server (the matured form of which today is
\href{https://github.com/punt-labs/postern}{Postern}), and tests them
by evaluating expressions.

\textbf{What Claude Code cannot do:} browse the Pharo image natively.
It has no \texttt{SystemBrowser} tool.  To compensate, the bootstrap
phase builds \cls{BootstrapImageBrowser} --- an introspection facade that
wraps Pharo's \texttt{SystemNavigation}, class hierarchy, and method
dictionaries behind text-returning methods callable via the eval server.

\textbf{Iteration loop:}
\begin{enumerate}
  \item Write \texttt{.st} file on disk
  \item POST \texttt{CodeImporter evaluateFileNamed:} to eval server
  \item POST test expression to eval server
  \item Read result, fix errors, repeat
\end{enumerate}

\textbf{Shared display:} The Pharo Transcript window is visible to
both the human (watching the Pharo GUI) and Claude Code (reading via
\texttt{Transcript contents} through the eval server).  This provides
real-time observability without additional infrastructure.

\subsection{Phase 2: External Agent Builds the Internal Agent}

\begin{center}
\begin{tabularx}{\linewidth}{@{}lY@{}}
\toprule
\textbf{Builder} & Claude Code (TypeScript CLI) \\
\textbf{Target} & \texttt{anthropic-sdk-pharo} \\
\textbf{Dependencies} & SDK from Phase~1 \\
\textbf{Key class} & \cls{PharoAgent} \\
\bottomrule
\end{tabularx}
\end{center}

Using the SDK built in Phase~1, Claude Code constructs the internal
agent runtime:

\begin{itemize}
  \item \cls{PharoAgent} --- the agent loop: receive spec, call Claude
    API (via the SDK), dispatch tool calls, evaluate code, iterate on
    errors.
  \item Tool definitions --- \texttt{browse\_class},
    \texttt{browse\_method}, \texttt{search\_code},
    \texttt{suggest\_selector}, \texttt{evaluate}, \texttt{lint}.
  \item Composable system prompt --- 6 class-side methods covering
    Smalltalk fundamentals, image model, Morphic reference, tool usage,
    and operational discipline.
  \item Instance-side \sel{mission:} for specialization.
\end{itemize}

\textbf{The critical difference from Phase~1:} the internal agent has
tools that the external agent does not.  \cls{PharoAgent} can browse
classes, read method source, fuzzy-match selectors, and evaluate code
--- all as direct method calls inside the image.  The external agent
must go through the eval server for every interaction.  The internal
agent operates at the speed of message sends.

\textbf{Validation:} the internal agent is tested by giving it UI
specs and verifying it produces working Morphic windows.  Spike~4
results: 3--5 turns for UI creation, modification, and cleanup.

\subsection{Phase 3: Internal Agent Builds Itself}

Phase 3 transferred self-improvement capability to the internal agent.
This work, along with the full agent runtime, continues in the separate
\texttt{claude-agent-sdk-pharo} repository.

\subsection{Phase 4: Internal Agent Extends Pharo --- By Building Anything}

\begin{center}
\begin{tabularx}{\linewidth}{@{}lY@{}}
\toprule
\textbf{Builder} & \cls{PharoAgent} (inside the image) \\
\textbf{Target} & Any in-image extension --- UI components, developer
  tools, integration services, packages \\
\textbf{Substrate} & Live Pharo image with full introspection \\
\textbf{Tools} & Postern, Iceberg, agent self-modification primitives \\
\bottomrule
\end{tabularx}
\end{center}

Once Phase~3 has handed the image over to the internal agent, Phase~4
is open-ended.  The agent does not need a specific external project to
target.  It can build new Morphic UI on demand without a fixed widget
vocabulary, add developer tools that extend the System Browser,
Inspector, or Test Runner, stand up integration services (HTTP
endpoints, MCP servers, message queues) that other systems consume,
define new packages that capture domain abstractions discovered during
use, and refactor and reorganize its own code as patterns emerge.

The point is not any one of these.  The external bootstrap was the
means; an in-image agent that can extend Pharo arbitrarily is the
result.  The Messaging SDK shipped here is the first such product.
Future work in \texttt{claude-agent-sdk-pharo} carries this capability
forward.

% ===================================================================
\section{Key Design Decisions}
% ===================================================================

\subsection{Why HTTP, Not MCP, for Bootstrap}

The bootstrap mechanism is a raw HTTP eval server, not an MCP server.

\begin{center}
\begin{tabularx}{\linewidth}{@{}lYY@{}}
\toprule
& \textbf{HTTP Eval} & \textbf{MCP Server} \\
\midrule
Setup       & 2 lines of Smalltalk (the minimal invocation) & JSON-RPC framing, tool schemas, transport \\
Latency     & Sub-millisecond eval & MCP protocol overhead \\
Flexibility & Arbitrary code       & Fixed tool definitions \\
Bootstrap   & Can build MCP server & Requires MCP server to exist \\
\bottomrule
\end{tabularx}
\end{center}

The eval server can build anything, including an MCP server.  An MCP
server cannot build itself.  The eval server is the minimal bootstrap;
MCP is the production interface built on top of it.

\subsection{Why Chunk Files, Not Tonel}

Pharo's modern packaging format is Tonel (one file per method, directory
per class).  The bootstrap uses chunk format (\texttt{.st} files with
\texttt{!} delimiters) because:

\begin{itemize}
  \item \texttt{CodeImporter evaluateFileNamed:} loads chunk files
    directly --- no Metacello, no Iceberg, no git integration.
  \item One file per class is manageable at bootstrap scale (~5 classes).
  \item Chunk format requires CR line endings (\texttt{0x0D}), which
    \texttt{.gitattributes} handles via \texttt{eol=cr}.
  \item Migration to Tonel/Metacello happens after the bootstrap is
    complete (Phase~3 or later).
\end{itemize}

\subsection{Why the Image is Disposable}

The Makefile enforces a critical invariant: \texttt{make rebuild} must
produce a working image from a fresh download plus the \texttt{.st}
files.  The image is convenient (faster startup, preserved state) but
never precious.  All code lives on disk.

This matters because:
\begin{itemize}
  \item Pharo images can corrupt (crash during save, stale socket state).
  \item \texttt{Smalltalk snapshot:andQuit:} captures running processes,
    open connections, and thread state --- all of which may be invalid
    on restore.
  \item The eval server and any network connections do not survive image
    save/restore reliably.
  \item A rebuild from source is the recovery path.
\end{itemize}

\subsection{Why the Transcript as Shared Display}

The Pharo Transcript serves as the coordination surface between the
external agent and the human:

\begin{itemize}
  \item The human sees it in the Pharo GUI.
  \item Claude Code reads it via
    \texttt{curl ... -d "Transcript contents"}.
  \item The internal agent logs tool calls and results there.
  \item No additional monitoring infrastructure is needed.
\end{itemize}

This is an emergent property of the eval server architecture: anything
visible in the image is readable by the external agent.  The Transcript
happens to be the simplest shared text surface.

% ===================================================================
\section{Bootstrap vs.\ Production Naming}
% ===================================================================

Bootstrap code must be clearly distinguished from production code to
avoid confusion when production classes arrive.

\subsection{Package Naming}

All bootstrap code uses the \pkg{Claude-Bootstrap} package prefix:

\begin{center}
\begin{tabularx}{\linewidth}{@{}llY@{}}
\toprule
\textbf{Bootstrap Package} & \textbf{Production Package} & \textbf{Contents} \\
\midrule
\pkg{Claude-Bootstrap-Server}  & evolved into
  \href{https://github.com/punt-labs/postern}{Postern}, a separate
  product &
  \cls{BootstrapEvalServer}: HTTP eval endpoint \\
\pkg{Claude-Bootstrap-Browser} & \pkg{Claude-Workbench} (in
  \texttt{claude-agent-sdk-pharo}) &
  \cls{BootstrapImageBrowser}: introspection facade \\
\pkg{Claude-Bootstrap-SDK}     & \pkg{Claude-Messaging-Client} &
  \cls{ClaudeClient}: spike API client \\
\pkg{Claude-Bootstrap-Agent}   & \pkg{Claude-Agent-Tools} (in
  \texttt{claude-agent-sdk-pharo}, not in this repo) &
  \cls{PharoAgent}: spike agent loop \\
\bottomrule
\end{tabularx}
\end{center}

\subsection{Transition Strategy}

When production classes arrive (built per the SDK specification),
they go into their production packages (\pkg{Claude-Messaging-Client},
\pkg{Claude-Agent-Tools}, etc.).  The bootstrap classes are then
removed:

\begin{enumerate}
  \item Production \cls{ClaudeClient} (from \pkg{Claude-Messaging-Client})
    replaces the bootstrap version.
  \item Production \cls{PharoAgent} (from \pkg{Claude-Agent-Tools})
    replaces the bootstrap version.
  \item \cls{BootstrapImageBrowser} migrates to \pkg{Claude-Workbench}
    with API refinements.
  \item \cls{BootstrapEvalServer} remains available for development but is
    not part of any production package --- it is scaffolding.
  \item The entire \pkg{Claude-Bootstrap} package tree can be removed
    from a production image.
\end{enumerate}

\textbf{Key invariant:} at no point should bootstrap and production
classes with the same name coexist in the image.  The transition is
a replacement, not an overlay.

% ===================================================================
\section{The Handoff Problem}
% ===================================================================

The most delicate moment in the bootstrap is the transition from
Phase~2 (external builds internal) to Phase~3 (internal builds itself).

\subsection{When to Hand Off}

The internal agent is ready for self-improvement when it can:

\begin{enumerate}
  \item Browse its own classes and methods accurately.
  \item Modify its own system prompt and observe the effect.
  \item Add a new tool to its registry and use it in the next
    conversation turn.
  \item Define a new class, compile methods on it, and verify with
    \sel{lint:}.
  \item Recover from errors (wrong selector, undeclared variable) using
    \sel{suggestSelector:onClass:} and retry.
\end{enumerate}

All five capabilities were demonstrated during the Spike~4 validation.

\subsection{What Cannot Be Handed Off}

Some responsibilities remain external as the internal agent matures:
git operations beyond Iceberg's commit/push/fetch (notably merge),
PR/CI/CD workflows on GitHub, and direct human communication via the
Claude Code terminal.

% ===================================================================
\section{Experimental Results}
% ===================================================================

\subsection{Spike 4: Agent Loop Validation}

The internal agent (\cls{PharoAgent}) was tested with 5 UI specs of
increasing complexity:

\begin{center}
\begin{tabular}{@{}lrl@{}}
\toprule
\textbf{Spec} & \textbf{Turns} & \textbf{Outcome} \\
\midrule
Simple window with label       & 1--2 & Success \\
Multi-column table (5 rows)    & 3--5 & Success (after prompt revision) \\
Dynamic info panel (live data) & 5    & Success \\
Button creation                & 2    & Success \\
World cleanup (delete morphs)  & 3    & Success \\
\bottomrule
\end{tabular}
\end{center}

\subsection{System Prompt Evolution}

The system prompt went through three versions:

\begin{center}
\begin{tabularx}{\linewidth}{@{}lrY@{}}
\toprule
\textbf{Version} & \textbf{Size} & \textbf{Result} \\
\midrule
v1 (minimal) & 0.5 KB & Agent guessed wrong APIs, hit max retries (20) \\
v2 (patterns) & 2 KB & Table test succeeded but inconsistently \\
v3 (comprehensive) & 4.5 KB & 3--5 turns reliably.  6 composable
  class-side methods. \\
\bottomrule
\end{tabularx}
\end{center}

\textbf{Key finding:} working code patterns in the system prompt are
more valuable than browsing tools alone.  The agent can discover APIs
through introspection, but discovering the right \emph{combination} of
API calls to achieve a goal is where pre-tested patterns save 10--15
turns.

\subsection{Streaming Implementation}

The SSE streaming implementation went through two approaches:

\begin{enumerate}
  \item \textbf{Real-time via ZnLineReader} --- blocked the Zinc handler
    thread because \texttt{nextLine} waits indefinitely after
    \texttt{message\_stop}.  Forking into a Process didn't help; the
    blocking read consumed the process.
  \item \textbf{Buffered via full body read} --- \texttt{ZnClient} reads
    the complete chunked response body as a string.  SSE events are
    parsed from the string.  No blocking reads.  Trade-off: no
    incremental delivery.
\end{enumerate}

The buffered approach remains available for callers who prefer it.
Real-time streaming via raw \texttt{SecureSocket} is implemented in
\cls{ClaudeStreamingSocket} (\pkg{Claude-Messaging-Streaming} package).

\subsection{Pharo 12 API Discoveries}

Several Pharo~12 API changes were discovered during the bootstrap:

\begin{center}
\begin{tabularx}{\linewidth}{@{}lY@{}}
\toprule
\textbf{Deprecated} & \textbf{Replacement} \\
\midrule
\texttt{cls organization categories} & \texttt{cls protocolNames} \\
\texttt{cls organization listAtCategoryNamed:} &
  \texttt{cls selectorsInProtocol:} \\
\texttt{ref actualClass} (on CompiledMethod) &
  \texttt{ref methodClass} \\
\texttt{FileStream fileIn:} & \texttt{CodeImporter evaluateFileNamed:} \\
\texttt{ZnBlockDelegate} & Does not exist in Pharo~12 \\
\texttt{Compiler} (unqualified) & \texttt{OpalCompiler} (or
  \texttt{self class compiler}) \\
\bottomrule
\end{tabularx}
\end{center}

These were all discovered by running code in the eval server and
reading the error messages or deprecation warnings in the Transcript.
The introspection toolkit (\cls{BootstrapImageBrowser}) was updated to use
the non-deprecated APIs, eliminating all warnings.

% ===================================================================
\section{Comparison with Alternative Approaches}
% ===================================================================

\subsection{Electron or Native UI}

A native UI framework (Electron, SwiftUI, Qt) would require:
\begin{itemize}
  \item A wire protocol between the agent and the UI process.
  \item A fixed widget vocabulary.
  \item No runtime introspection or code modification.
  \item Separate build/compile/deploy cycles.
\end{itemize}

Pharo eliminates all four by making the UI framework and the agent
runtime the same process with shared objects.

% ===================================================================
\section{Security Considerations}
% ===================================================================

The HTTP eval server executes arbitrary Smalltalk code from any process
that can reach port~8422 on localhost.  This is by design --- the
bootstrap requires unrestricted code execution.  Security boundaries:

\begin{itemize}
  \item The server listens on \texttt{localhost} only (not
    \texttt{0.0.0.0}).
  \item No authentication --- any local process can evaluate code.
  \item The image runs with the user's OS permissions.
  \item The eval server should not be exposed to a network.
  \item For production, the eval server is replaced by a structured MCP
    interface with defined tools (not arbitrary eval).
\end{itemize}

The eval server is a development tool, not a production service.  It is development
scaffolding that is replaced by a structured MCP interface for production use.

The mature version of this eval server,
\href{https://github.com/punt-labs/postern}{Postern}, retains this
localhost-only model with an enhanced \texttt{/help} endpoint and
per-port isolation.

% ===================================================================
\section{Lessons Learned}
% ===================================================================

\begin{enumerate}[label=\textbf{L\arabic*}.]
  \item \textbf{The simplest bridge wins.}  The HTTP eval server is 2
    lines of Smalltalk.  Every more sophisticated approach (MCP, file
    watchers, socket protocols) requires the eval server to bootstrap
    itself.  Start with eval; build everything else on top.

  \item \textbf{Working patterns beat browsing tools.}  The agent can
    browse the entire Morphic class hierarchy, but discovering the right
    API combination for ``show a 3-column table'' takes 20+ turns.  A
    tested code pattern in the system prompt reduces this to 3 turns.
    Browsing tools are for error recovery and edge cases, not primary
    navigation.

  \item \textbf{The Transcript is a natural coordination surface.}  In
    hindsight, any object visible in the image is a potential
    coordination point.  The Transcript is the obvious one because it is
    append-only text, always visible, and readable via the eval server.

  \item \textbf{Image save is dangerous during development.}  Saving
    with open connections, running processes, or stale eval server state
    produces images that fail on restart.  The rebuild-from-source
    discipline (\texttt{make rebuild}) is essential, not optional.

  \item \textbf{CR line endings are a real obstacle.}  Pharo's chunk
    format requires \texttt{0x0D} line endings.  Every tool in the Unix
    ecosystem produces \texttt{0x0A}.  This caused ``Method source
    contains linefeeds'' warnings on every file-in until we converted
    all files and added \texttt{.gitattributes} rules.

  \item \textbf{The LLM's Smalltalk is better than expected.}  With
    the comprehensive system prompt (Smalltalk fundamentals, Morphic
    reference, Pharo~12 specifics), the model produces correct Smalltalk
    in 2--3 attempts.  The main failure mode is using deprecated APIs
    from older Pharo versions, which the error hints and
    \sel{suggestSelector:onClass:} tool resolve.

  \item \textbf{Pharo's cooperative scheduling matters.}  Forked
    processes do not preempt each other at the same priority.  The
    streaming SSE reader ``blocked'' because it never yielded to the UI
    thread.  The fix was not better threading but avoiding blocking
    reads entirely (buffered body approach).

  \item \textbf{The bootstrap is the product.}  The eval server became
    \href{https://github.com/punt-labs/postern}{Postern}, a standalone
    designed product.  The introspection toolkit
    (\cls{BootstrapImageBrowser}) migrates to the Workbench layer (in
    \texttt{claude-agent-sdk-pharo}).  The composable system prompt is
    the foundation of the internal agent.

  \item \textbf{Self-improvement is a spectrum, not a threshold.}
    The internal agent can modify its own classes after Phase~2, but it
    still needs the external agent for git, testing discipline, and
    project management.  Full autonomy is not the goal; effective
    collaboration between two agents with different strengths is.
\end{enumerate}

% ===================================================================
\section{Future Work}
% ===================================================================

\begin{enumerate}
  \item \textbf{MCP server in Pharo} (Spike~2) --- replace the eval
    server with a structured MCP interface for production use; tracked
    in \texttt{claude-agent-sdk-pharo}.
  \item \textbf{Claude Code to Pharo bridge} (Spike~5) --- full
    two-agent round-trip: spec in, UI out, events back; tracked in
    \texttt{claude-agent-sdk-pharo}.
  \item \textbf{Claude Workbench} --- Morphic application replicating
    the Anthropic Console Workbench, serving as both SDK demo and
    development tool; tracked in \texttt{claude-agent-sdk-pharo}.
  \item \textbf{Self-improvement validation} --- systematic testing of
    the internal agent's ability to modify its own code, add tools, and
    extend the system prompt; tracked in
    \texttt{claude-agent-sdk-pharo}.
  \item \textbf{Research report} --- formal study of the bootstrap
    process, comparing development velocity, error rates, and turn
    counts across the four phases.
\end{enumerate}

\end{document}
